% ============================================================
%  CHAPTER 3 — System Architecture
% ============================================================

\chapter{System Architecture}
\label{ch:architecture}

\section{Global Architecture Diagram}

Figure~\ref{fig:global_arch} presents the global architecture of the complete system,
showing all five layers defined in the project specification: the sensor layer, the
embedded layer, the server layer, the AI layer, and the application layer.

\begin{figure}[H]
	\centering
	\includegraphics[width=\linewidth]{global_architecture_diagram.png}
	\caption{Global architecture of the system.}
	\label{fig:global_arch}
\end{figure}

% flowchart LR
%     subgraph SL["Sensor layer"]
%         DHT["DHT22\nTemp / Hum"]
%         MAX["MAX30102\nHR / SpO₂"]
%     end
%
%     subgraph ESP["Embedded layer (ESP32)"]
%         C1["Core 1\nSensors + web server"]
%         C0["Core 0\nFirebase task"]
%         AP["Wi-Fi AP\n192.168.4.1"]
%         STA["Wi-Fi STA\n802.11n"]
%         FLASH["LittleFS + NVS\nFlash storage"]
%     end
%
%     subgraph SRV["Server layer"]
%         AUTH["Firebase Auth\nemail / password"]
%         RTDB["Firebase RTDB\n/rooms  /patients"]
%     end
%
%     subgraph AIL["AI layer"]
%         AI["AI / Cloud\nPredictions"]
%     end
%
%     subgraph APP["Application layer"]
%         WEB["Web dashboard"]
%         CFG["Config browser"]
%     end
%
%     DHT -->|data| C1
%     MAX -->|data| C1
%     C1 <-->|shared globals| C0
%     C1 --> AP
%     C0 --> STA
%     FLASH -.->|NVS load| C0
%     C1 -.-> FLASH
%     STA -->|auth| AUTH
%     STA -->|push| RTDB
%     RTDB --> AI
%     RTDB -->|read| WEB
%     AI -->|reports| WEB
%     AP <-->|HTTP config| CFG

\section{Functional Architecture}

The system is functionally decomposed into four independent subsystems, each with a
well-defined interface:

\begin{description}
	\item[Sensor Subsystem (ESP32, Core~1).]
	\item[Network/Firebase Subsystem (ESP32, Core~0).]
	\item[Settings/Web Server Subsystem (ESP32, Core~1).]
	\item[Configuration and Globals Layer.]
\end{description}

% -------------------------------------------------------
% TODO (Backend Team): Describe the functional decomposition
% of the backend (API router, auth middleware, alert engine,
% database access layer, etc.).
% -------------------------------------------------------
\begin{center}
	\textit{[Backend functional decomposition to be added.]}
\end{center}

% -------------------------------------------------------
% TODO (Web / AI Teams): Describe the functional
% decomposition of the web dashboard and AI subsystem.
% -------------------------------------------------------
\begin{center}
	\textit{[Web dashboard and AI functional decomposition to be added.]}
\end{center}

\section{Hardware Architecture}

\begin{description}
	\item[Microcontroller.]
	      ESP32-WROOM-32, dual-core at 240\,MHz,
	      520\,KiB SRAM, 4\,MiB flash.

	\item[Temperature and Humidity Sensor.]
	      DHT22, single-bus digital interface on GPIO~4.

	\item[Pulse Oximeter and Heart Rate Sensor.]
	      MAX30102, I\textsuperscript{2}C fast mode (400\,kHz), SDA: GPIO~21,
	      SCL: GPIO~22.

	\item[Communication.]
	      IEEE 802.11b/g/n Wi-Fi radio integrated on-chip.
	      AP interface (192.168.4.1) serves the configuration portal.
\end{description}

% -------------------------------------------------------
% TODO (Backend Team): Describe the server hardware
% (cloud VM, hosting tier, database engine, etc.).
% -------------------------------------------------------
\begin{center}
	\textit{[Server hardware description to be added by the backend team.]}
\end{center}

\section{Software Architecture}

\subsection*{Embedded Software (ESP32)}

The firmware is organised into six source files following a strict, acyclic include
hierarchy.
Table~\ref{tab:sw_modules} summarises the modules, their execution core, and their
responsibilities.

\begin{table}[H]
	\centering
	\caption{ESP32 firmware module summary.}
	\label{tab:sw_modules}
	\begin{tabularx}{\linewidth}{lllX}
		\toprule
		\textbf{File}           & \textbf{Context} & \textbf{Responsibility}                                                  \\
		\midrule
		\texttt{globals.h}      & Header           & Shared variable declarations, DBG macros, timing constants               \\
		\texttt{sensors.h/cpp}  & Loop             & DHT22 and MAX30102 acquisition, BPM and SpO\textsubscript{2} computation \\
		\texttt{network.h/cpp}  & FreeRTOS task    & Firebase auth, upload, Wi-Fi management                                  \\
		\texttt{settings.h/cpp} & ISR callbacks    & HTTP routes, LittleFS, NVS persistence                                   \\
		\texttt{esp32.ino}      & Entry point      & Global object definitions, \texttt{setup()}, \texttt{loop()}             \\
		\bottomrule
	\end{tabularx}
\end{table}

The dependency graph is strictly acyclic: \texttt{esp32.ino} includes all modules;
\texttt{network.h/cpp} and \texttt{settings.h/cpp} include \texttt{globals.h};
\texttt{sensors.h/cpp} includes \texttt{globals.h}.
No module includes another module's header directly.

% -------------------------------------------------------
% TODO (Backend Team): Describe the backend software
% architecture (framework, router structure, middleware
% chain, database ORM or driver, deployment environment).
% -------------------------------------------------------
\subsection*{Backend Software}
\begin{center}
	\textit{[To be completed by the backend team.]}
\end{center}

% -------------------------------------------------------
% TODO (Web Dashboard Team): Describe the web application
% architecture (framework, component tree, state management,
% API client, chart libraries, etc.).
% -------------------------------------------------------
\subsection*{Web Dashboard Software}
\begin{center}
	\textit{[To be completed by the web dashboard team.]}
\end{center}

% -------------------------------------------------------
% TODO (AI Team): Describe the AI software architecture
% (framework, model architecture, training pipeline,
% deployment mechanism — cloud API, ONNX, TFLite, etc.).
% -------------------------------------------------------
\subsection*{AI Software}
\begin{center}
	\textit{[To be completed by the AI team.]}
\end{center}

\section{Processing Pipeline}

The end-to-end data processing pipeline for a single sensor reading, from hardware
to cloud record, proceeds as follows:


% -------------------------------------------------------
% TODO (Backend / AI / Web Teams): Extend steps 9 onward
% with the full downstream processing pipeline (backend
% alert evaluation, AI inference, dashboard update cycle).
% -------------------------------------------------------
\begin{center}
	\textit{[Downstream pipeline details to be added by the backend, AI, and web teams.]}
\end{center}
